\documentclass[tikz,border=8pt]{standalone}

\usepackage{amsmath}
\usepackage{bm}
\usepackage{newtxtext,newtxmath}
\usepackage{tikz-3dplot}
\usetikzlibrary{arrows.meta,calc}

\begin{document}

% ------------------------------------------------------------
% 3D view
% ------------------------------------------------------------
\tdplotsetmaincoords{65}{115}

\begin{tikzpicture}[
    tdplot_main_coords,
    line cap=round,
    line join=round,
    >=Stealth,
    font=\small
]

% ------------------------------------------------------------
% Parameters
% ------------------------------------------------------------
\def\R{2.45}

% tangent point on the hemisphere
\def\thp{43}    % polar angle from +z
\def\php{322}   % azimuth angle

% size of the small tangent patch
\def\su{0.38}
\def\sv{0.30}

% ------------------------------------------------------------
% Styles
% ------------------------------------------------------------
\tikzset{
    spheregrid/.style={
        draw=gray!32,
        line width=.55pt
    },
    sphereedge/.style={
        draw=gray!58,
        line width=1.05pt
    },
    hiddenedge/.style={
        draw=gray!35,
        dashed,
        dash pattern=on 2pt off 2.4pt,
        line width=.65pt
    },
    tangentpatch/.style={
        draw=black!60,
        line width=.8pt,
        fill=black!35,
        fill opacity=.16
    },
    ray/.style={
        draw=black!55,
        dashed,
        dash pattern=on 2.6pt off 2.4pt,
        line width=.75pt
    },
    volumeedge/.style={
        draw=black!50,
        line width=.55pt,
        fill=black!20,
        fill opacity=.045
    },
    point/.style={
        circle,
        fill=black!70,
        inner sep=1.05pt
    },
    centerpoint/.style={
        circle,
        fill=black!75,
        inner sep=1.15pt
    },
    labelbox/.style={
        fill=white,
        fill opacity=.86,
        text opacity=1,
        rounded corners=1.7pt,
        inner xsep=3pt,
        inner ysep=1.6pt,
        align=center
    }
}

% ------------------------------------------------------------
% Draw an upper hemisphere z >= 0
% ------------------------------------------------------------
\newcommand{\drawhemi}{%
    % equator: back dashed, front solid
    \draw[hiddenedge]
        plot[domain=25:205, samples=90, variable=\t]
        ({\R*cos(\t)}, {\R*sin(\t)}, 0);

    \draw[sphereedge]
        plot[domain=205:385, samples=90, variable=\t]
        ({\R*cos(\t)}, {\R*sin(\t)}, 0);

    % latitude circles
    \foreach \th in {25,48,70}{
        \pgfmathsetmacro{\rr}{\R*sin(\th)}
        \pgfmathsetmacro{\zz}{\R*cos(\th)}
        \draw[spheregrid]
            plot[domain=25:385, samples=110, variable=\t]
            ({\rr*cos(\t)}, {\rr*sin(\t)}, \zz);
    }

    % meridians
    \foreach \ph in {25,70,115,160,205,250,295,340}{
        \draw[spheregrid]
            plot[domain=0:90, samples=80, variable=\t]
            ({\R*sin(\t)*cos(\ph)},
             {\R*sin(\t)*sin(\ph)},
             {\R*cos(\t)});
    }

    % two stronger silhouette-like meridians
    \draw[sphereedge]
        plot[domain=0:90, samples=90, variable=\t]
        ({\R*sin(\t)*cos(25)},
         {\R*sin(\t)*sin(25)},
         {\R*cos(\t)});

    \draw[sphereedge]
        plot[domain=0:90, samples=90, variable=\t]
        ({\R*sin(\t)*cos(205)},
         {\R*sin(\t)*sin(205)},
         {\R*cos(\t)});
}

% ------------------------------------------------------------
% Main point on the hemisphere
% ------------------------------------------------------------
\coordinate (O) at (0,0,0);

\pgfmathsetmacro{\px}{\R*sin(\thp)*cos(\php)}
\pgfmathsetmacro{\py}{\R*sin(\thp)*sin(\php)}
\pgfmathsetmacro{\pz}{\R*cos(\thp)}

\coordinate (P) at (\px,\py,\pz);

% ------------------------------------------------------------
% Local tangent basis at P on the sphere
%
% n = P / R
% u = normalized dX / dphi
% v = normalized dX / dtheta
% ------------------------------------------------------------
\pgfmathsetmacro{\nx}{sin(\thp)*cos(\php)}
\pgfmathsetmacro{\ny}{sin(\thp)*sin(\php)}
\pgfmathsetmacro{\nz}{cos(\thp)}

% u direction, tangent along azimuth
\pgfmathsetmacro{\uxraw}{-sin(\php)}
\pgfmathsetmacro{\uyraw}{ cos(\php)}
\pgfmathsetmacro{\uzraw}{0}

% v direction, tangent along polar angle
\pgfmathsetmacro{\vxraw}{cos(\thp)*cos(\php)}
\pgfmathsetmacro{\vyraw}{cos(\thp)*sin(\php)}
\pgfmathsetmacro{\vzraw}{-sin(\thp)}

% these are already unit length for sphere parameterization
\pgfmathsetmacro{\ux}{\uxraw}
\pgfmathsetmacro{\uy}{\uyraw}
\pgfmathsetmacro{\uz}{\uzraw}

\pgfmathsetmacro{\vx}{\vxraw}
\pgfmathsetmacro{\vy}{\vyraw}
\pgfmathsetmacro{\vz}{\vzraw}

% ------------------------------------------------------------
% Four vertices of the small tangent patch
% ------------------------------------------------------------
\coordinate (Q1) at ({\px+\su*\ux+\sv*\vx},
                     {\py+\su*\uy+\sv*\vy},
                     {\pz+\su*\uz+\sv*\vz});

\coordinate (Q2) at ({\px-\su*\ux+\sv*\vx},
                     {\py-\su*\uy+\sv*\vy},
                     {\pz-\su*\uz+\sv*\vz});

\coordinate (Q3) at ({\px-\su*\ux-\sv*\vx},
                     {\py-\su*\uy-\sv*\vy},
                     {\pz-\su*\uz-\sv*\vz});

\coordinate (Q4) at ({\px+\su*\ux-\sv*\vx},
                     {\py+\su*\uy-\sv*\vy},
                     {\pz+\su*\uz-\sv*\vz});

% ------------------------------------------------------------
% Draw order
% ------------------------------------------------------------
\drawhemi

% subtle volume faces
\filldraw[volumeedge] (O)--(Q1)--(Q2)--cycle;
\filldraw[volumeedge] (O)--(Q2)--(Q3)--cycle;
\filldraw[volumeedge] (O)--(Q3)--(Q4)--cycle;
\filldraw[volumeedge] (O)--(Q4)--(Q1)--cycle;

% dashed rays from center to the four vertices
\draw[ray] (O) -- (Q1);
\draw[ray] (O) -- (Q2);
\draw[ray] (O) -- (Q3);
\draw[ray] (O) -- (Q4);

% tangent patch
\filldraw[tangentpatch] (Q1)--(Q2)--(Q3)--(Q4)--cycle;

% vertices and center
\node[centerpoint] at (O) {};
\node[point] at (Q1) {};
\node[point] at (Q2) {};
\node[point] at (Q3) {};
\node[point] at (Q4) {};

% ------------------------------------------------------------
% dV label inside the cone / pyramid volume
% ------------------------------------------------------------
\coordinate (Vlabel) at ({0.42*\px+0.18},
                         {0.42*\py+0.05},
                         {0.42*\pz});

\node[labelbox, text=black] at (Vlabel) {$\mathrm{d}V$};

\end{tikzpicture}

\end{document}